\documentclass[UTF8,11pt]{ctexart}

\usepackage[a4paper,margin=22mm]{geometry}
\usepackage{amsmath,amssymb,mathtools}
\usepackage{booktabs,array,longtable}
\usepackage{enumitem}
\usepackage{xcolor}
\usepackage{hyperref}
\usepackage{tikz}
\usetikzlibrary{calc,intersections}

\hypersetup{colorlinks=true,linkcolor=blue!55!black,urlcolor=blue!65!black}
\setlength{\parindent}{2em}
\setlength{\parskip}{0.35em}
\setlist[itemize]{leftmargin=2em,itemsep=0.2em}
\setlist[enumerate]{leftmargin=2.2em,itemsep=0.25em}

\newcommand{\dirseg}[2]{\overline{#1#2}}
\newcommand{\polar}{\operatorname{polar}}

\title{射影几何与极点极线理论（上）\\视频内容提取与数学化整理}
\author{}
\date{整理日期：2026年8月5日}

\begin{document}
\maketitle

\section*{来源与整理原则}

\noindent 视频来源：Tutor王博士数学，
\href{https://www.bilibili.com/video/BV1Me411o7ku}
{《射影几何与极点极线理论（上）》}，时长约 $69$ 分钟。

本文依据完整视频音频、板书画面和本地语音转写整理。它不是逐字稿，而是按照视频实际讲授顺序恢复概念、推导和例题。口语中的“交笔、调和线书、极限”等识别错误，分别校正为“交比、调和线束、极线”。

视频的中心意图不是让学生生搬硬套极线方程，而是建立如下来源链：
\[
\boxed{\text{射影不变性}}
\Longrightarrow
\boxed{\text{交比}}
\Longrightarrow
\boxed{\text{调和点列与调和线束}}
\Longrightarrow
\boxed{\text{极点--极线}}
\Longrightarrow
\boxed{\text{自极三角形与定点问题}}.
\]

\section{视频时间索引}

\begin{longtable}{>{\ttfamily\small}p{2.7cm}p{11.0cm}}
\toprule
时间 & 内容 \\
\midrule
00:00--02:57 & 为什么从射影几何理解圆锥曲线命题；射影几何研究什么在投影下保持不变。\\
02:57--08:03 & 四点交比的定义；利用面积比、正弦比说明交比的射影不变性。\\
08:03--13:25 & 交比为 $-1$；调和比、调和点列、调和分割和调和共轭。\\
13:25--17:51 & 调和平均关系；调和点列仍具有射影不变性。\\
17:51--23:08 & 完全四边形产生调和点列；视频指出可由 Menelaus 定理和 Ceva 定理证明。\\
23:08--28:44 & 调和线束；把调和关系写成四条直线斜率之间的关系。\\
28:44--33:31 & 调和线束的中点模型；一条截线平行于线束中某条射线时，调和关系化为中点关系。\\
33:31--40:00 & 从圆锥曲线割线上的调和共轭点出发，用定比分点和点差法推出极线方程。\\
40:00--45:14 & 极线与切点弦方程是同一对象；数值例子 $x^2/4+y^2/2=1$，$P=(4,1)$。\\
45:14--47:08 & 极点极线的互易性，即 La Hire 定理的解析来源。\\
47:08--53:26 & 圆锥曲线内接完全四边形；三个对角点组成自极三角形。\\
53:26--55:27 & 用自极三角形解释高考中的定点、定直线结构。\\
55:27--59:37 & 焦点与准线互为极点极线；调和平均在圆锥曲线中的体现。\\
59:37--62:47 & 学习方法：按结构给高考题分类，而不是逐题积累孤立技巧。\\
62:47--69:04 & 椭圆 $x^2/9+y^2/4=1$ 的定点例题；由自极三角形转化为动点极线过定点。\\
\bottomrule
\end{longtable}

\section{射影不变性与交比}

\subsection{视频的出发点}

投影一般不保持长度、角度和普通线段比。射影几何关心的是：图形经过中心投影之后，哪些关系仍然保持不变？视频首先引入的射影不变量就是四点交比。

设 $A,C,B,D$ 为同一直线上的四点，所有线段均取有向线段。按照视频采用的次序，定义
\begin{equation}
\label{eq:cross-ratio}
(A,B;C,D)
=\frac{\dirseg AC/\dirseg AD}{\dirseg BC/\dirseg BD}
=\frac{\dirseg AC\,\dirseg BD}{\dirseg AD\,\dirseg BC}.
\end{equation}
不同教材对四个字母的排列可能不同，使用时必须先核对定义，不能只看“交比”两个字便套公式。

\subsection{为什么交比在中心投影下不变}

从直线外一点 $O$ 向 $A,C,B,D$ 引四条射线。由三角形面积公式，
\[
\frac{\dirseg AC}{\dirseg AD}
=\frac{S_{\triangle AOC}}{S_{\triangle AOD}},
\qquad
\frac{\dirseg BC}{\dirseg BD}
=\frac{S_{\triangle BOC}}{S_{\triangle BOD}}.
\]
进一步得到
\begin{equation}
\label{eq:sine-cross-ratio}
(A,B;C,D)
=\frac{\sin\angle AOC}{\sin\angle AOD}
 \frac{\sin\angle BOD}{\sin\angle BOC}.
\end{equation}
右端只依赖四条射线的方向，不依赖用哪一条截线截取它们。因此，只要投影中心和四条投影射线不变，换一条截线得到的新四点具有相同交比。

这正是视频所谓“交比具有射影不变性”的初等解释。

\section{调和点列、调和分割与调和平均}

\subsection{调和点列}

当
\begin{equation}
\label{eq:harmonic}
(A,B;C,D)=-1
\end{equation}
时，称 $A,C,B,D$ 构成一组\textbf{调和点列}，也称 $C,D$ 调和分割 $AB$；$C,D$ 互为关于点对 $A,B$ 的调和共轭点。

视频用参数 $\lambda>0$（$\lambda\ne1$）表示为
\[
\overrightarrow{AC}=\lambda\overrightarrow{CB},
\qquad
\overrightarrow{AD}=-\lambda\overrightarrow{DB}.
\]
一个点内分 $AB$，另一个点按同一比值外分 $AB$，方向差异使交比为 $-1$。

调和关系具有互易性：若 $C,D$ 调和分割 $A,B$，那么换一种排列后，$A,B$ 也调和分割 $C,D$。这为后面极点极线的互易性作了铺垫。

\subsection{调和平均关系}

在视频所画的点序下，由调和条件可整理为
\begin{equation}
\label{eq:harmonic-mean}
\boxed{
\frac{1}{\dirseg AC}+\frac{1}{\dirseg AD}
=\frac{2}{\dirseg AB}
}.
\end{equation}
注意式中仍是有向线段。它说明 $\dirseg AB$ 在倒数意义下是 $\dirseg AC$ 与 $\dirseg AD$ 的调和平均，这也是“调和”一词在视频中的数量解释。

\section{完全四边形：制造调和关系}

取四个点 $A,B,C,D$，其中任意三点不共线，连接六条边线便得到完全四边形（也常称完全四点形）。其三组对边交点称为三个对角点：
\[
P=AB\cap CD,qquad
M=AC\cap BD,qquad
N=AD\cap BC.
\]

视频在约 $17{:}51$ 处引入这一结构，并通过延长边、连接对角点，得到多组调和点列。其证明在视频中没有完整展开，只说明可结合 Menelaus 定理和 Ceva 定理完成。

这里真正需要掌握的不是某一幅固定字母图，而是：
\begin{enumerate}
\item 完全四边形能够产生调和点列；
\item 调和点列可以通过中心投影搬到另一条直线上；
\item 当 $A,B,C,D$ 位于同一圆锥曲线上时，三个对角点 $P,M,N$ 进一步组成自极三角形，见第~\ref{sec:self-polar} 节。
\end{enumerate}

\section{调和线束及其斜率形式}

若从一点 $O$ 发出的四条直线 $OA,OB,OC,OD$ 被任意截线截得的四点构成调和点列，则称这四条直线构成\textbf{调和线束}。由于交比具有射影不变性，这一定义与所选截线无关。

取一条水平截线，设 $O$ 到截线的距离为 $h$，四条直线的斜率依次为 $k_1,k_2,k_3,k_4$。若垂足为 $H$，则
\[
\dirseg AH=\frac{h}{k_1},\quad
\dirseg BH=\frac{h}{k_2},\quad
\dirseg CH=\frac{h}{k_3},\quad
\dirseg DH=\frac{h}{k_4}.
\]
把它们代入调和条件，公共因子约去后得到视频重点强调的斜率关系：
\begin{equation}
\label{eq:harmonic-slopes}
\boxed{
\frac{k_2-k_1}{k_4-k_1}
=\frac{k_3-k_2}{k_4-k_3}
}.
\end{equation}

这一形式把射影几何中的调和线束，直接翻译成高中解析几何可处理的斜率关系。若有竖直线，应改用方向参数或先旋转坐标，不能把“斜率不存在”直接代入式~\eqref{eq:harmonic-slopes}。

\subsection{中点模型}

若一条截线平行于调和线束中的某条射线，那么该射线与截线的交点可理解为无穷远点。调和点列中一个点移到无穷远后，交比为 $-1$ 的条件退化为中点关系。

视频给出的典型图式是：一条截线经过某一点 $A$，并平行于线束中的一条射线；它与另外两条射线交于 $M,N$，则
\[
\boxed{AM=AN}.
\]
视频用两次三角形相似证明这一结论，并指出许多高考题会把调和结构隐藏为“平行线 $+$ 中点”。反过来，看到中点和平行关系，也可尝试把平行方向补成无穷远点，恢复调和线束。

\section{从调和共轭点推导极线}

考虑椭圆
\begin{equation}
\label{eq:ellipse}
\mathcal E:\quad \frac{x^2}{a^2}+\frac{y^2}{b^2}=1.
\end{equation}
固定一点 $P(x_P,y_P)$，任作一条过 $P$ 的割线，交椭圆于
$A(x_1,y_1)$、$B(x_2,y_2)$。在割线上取点 $Q(x_Q,y_Q)$，使 $P,Q$ 关于 $A,B$ 调和共轭。

按视频的定比分点写法，存在参数 $\lambda$，使
\[
P=\frac{A-\lambda B}{1-\lambda},
\qquad
Q=\frac{A+\lambda B}{1+\lambda}.
\]
于是
\[
x_P=\frac{x_1-\lambda x_2}{1-\lambda},\qquad
x_Q=\frac{x_1+\lambda x_2}{1+\lambda},
\]
并有对应的 $y$ 坐标表达式。

因为 $A,B\in\mathcal E$，
\[
\frac{x_1^2}{a^2}+\frac{y_1^2}{b^2}=1,
\qquad
\frac{x_2^2}{a^2}+\frac{y_2^2}{b^2}=1.
\]
将第二式乘以 $\lambda^2$ 后与第一式相减，利用平方差分解和上述定比分点表达式，可得
\begin{equation}
\label{eq:polar-ellipse}
\boxed{
\frac{x_Px_Q}{a^2}+\frac{y_Py_Q}{b^2}=1
}.
\end{equation}

因此，无论割线怎样绕 $P$ 转动，与 $P$ 调和共轭的点 $Q$ 总在固定直线
\begin{equation}
\label{eq:polar-line}
\boxed{
\ell_P:\quad \frac{x_Px}{a^2}+\frac{y_Py}{b^2}=1
}
\end{equation}
上。这条直线称为 $P$ 关于椭圆的\textbf{极线}，$P$ 称为 $\ell_P$ 的\textbf{极点}。

这个推导是视频最重要的内容：
\[
\boxed{
\text{极线是过 }P\text{ 的各条割线上，}P\text{ 的调和共轭点的轨迹}
}.
\]
因此，“把 $x^2$ 换成 $x_Px$”不是无来源的替换规则，而是调和关系经过定比分点和点差法后的结果。

\subsection{极线的三种面貌}

对高中圆锥曲线题，可以把同一条极线理解为：
\begin{enumerate}
\item \textbf{调和面貌：}割线 $PAB$ 与极线交于 $Q$，则 $P,Q$ 关于 $A,B$ 调和共轭；
\item \textbf{切线面貌：}若 $P$ 在曲线外，从 $P$ 引两条切线，切点为 $T_1,T_2$，则 $T_1T_2$ 就是 $P$ 的极线；
\item \textbf{解析面貌：}标准方程通过双线性化得到极线方程。
\end{enumerate}

常见标准式为
\[
\begin{array}{c|c}
\text{圆锥曲线} & P(x_0,y_0)\text{ 的极线}\\ \hline
\dfrac{x^2}{a^2}+\dfrac{y^2}{b^2}=1
& \dfrac{x_0x}{a^2}+\dfrac{y_0y}{b^2}=1\\[6pt]
\dfrac{x^2}{a^2}-\dfrac{y^2}{b^2}=1
& \dfrac{x_0x}{a^2}-\dfrac{y_0y}{b^2}=1\\[6pt]
y^2=2px & yy_0=p(x+x_0)
\end{array}
\]

\subsection{视频中的数值例子}

对椭圆
\[
\frac{x^2}{4}+\frac{y^2}{2}=1
\]
和点 $P=(4,1)$，极线为
\[
\frac{4x}{4}+\frac{y}{2}=1,
\qquad\text{即}\qquad
\boxed{x+\frac{y}{2}=1}.
\]
视频借此说明：高考题若给出割线上的比例关系，要求证明某点 $Q$ 在定直线上，本质上可能就是在要求证明 $Q$ 位于 $P$ 的极线上。

\section{极点极线的互易性}

由式~\eqref{eq:polar-ellipse} 的对称性，
\[
Q\in\ell_P
\iff
\frac{x_Px_Q}{a^2}+\frac{y_Py_Q}{b^2}=1
\iff
P\in\ell_Q.
\]
因此
\begin{equation}
\label{eq:lahire}
\boxed{Q\in\polar(P)\iff P\in\polar(Q)}.
\end{equation}
这就是极点极线的互易性，通常称为 \textbf{La Hire 定理}。

它带来一个极其重要的定点工具：
\begin{quote}
若动点 $Q$ 始终在固定直线 $l$ 上，则 $Q$ 的极线始终经过直线 $l$ 的极点。
\end{quote}
也就是说，\textbf{点在定直线上运动}与\textbf{它的极线经过定点}互为对偶。

\section{完全四边形与自极三角形}
\label{sec:self-polar}

设 $A,B,C,D$ 是同一非退化圆锥曲线上的四点，三个对角点为
\[
P=AB\cap CD,qquad
M=AC\cap BD,qquad
N=AD\cap BC.
\]
视频把完全四边形产生的调和点列与极线定义结合，得到
\begin{equation}
\label{eq:self-polar}
\boxed{
\polar(P)=MN,qquad
\polar(M)=PN,qquad
\polar(N)=PM
}.
\end{equation}
因此三角形 $PMN$ 的每一边都是对面顶点的极线，称为关于该圆锥曲线的\textbf{自极三角形}。

这一结论解释了许多题中看似偶然的共线和定点：当四个圆锥曲线上的点构成完全四边形时，三组对边交点并不是三个无关的点，它们已经被极点极线关系锁定。

\section{焦点与准线的极点极线关系}

对椭圆
\[
\frac{x^2}{a^2}+\frac{y^2}{b^2}=1,qquad c^2=a^2-b^2,
\]
右焦点为 $F(c,0)$。由极线公式，$F$ 的极线为
\[
\frac{cx}{a^2}=1,
\qquad\text{即}\qquad
\boxed{x=\frac{a^2}{c}},
\]
它恰好是椭圆的右准线。因此，焦点与对应准线互为极点和极线。

这使“焦点、准线、割线、调和点列”统一起来：准线上的点的极线都经过对应焦点；反过来，经过焦点的直线的极点位于对应准线上。

视频进一步指出，一些焦半径倒数关系可以从调和平均的角度理解。这种联系适合作为结论来源，但在高中书面解答中仍应根据题设写成距离、坐标或韦达关系，不能只写“由调和性显然”。

\section{视频末尾例题：定点（3,2）从哪里来}

考虑椭圆
\begin{equation}
\label{eq:example-ellipse}
\frac{x^2}{9}+\frac{y^2}{4}=1.
\end{equation}
记四个顶点为
\[
A=(0,2),\quad B=(-3,0),\quad C=(0,-2),\quad D=(3,0).
\]
取椭圆上异于顶点的动点 $P$。按视频图形，定义
\[
M=AP\cap CD,qquad
N=DP\cap AC,qquad
Q=PC\cap AD.
\]
四点 $A,D,C,P$ 构成内接完全四边形，$M,N,Q$ 是它的三个对角点。因此由自极三角形结论，
\[
\boxed{MN=\polar(Q)}.
\]

另一方面，$Q$ 始终位于固定直线 $AD$ 上。由
$A=(0,2)$、$D=(3,0)$ 得
\begin{equation}
\label{eq:AD}
AD:\quad y=-\frac23x+2.
\end{equation}
设
\[
Q=\left(t,-\frac23t+2\right).
\]
点 $Q$ 关于椭圆~\eqref{eq:example-ellipse} 的极线，即 $MN$，为
\[
\frac{tx}{9}
+\frac{\left(-\frac23t+2\right)y}{4}=1.
\]
整理为
\begin{equation}
\label{eq:family-polar}
t\left(\frac{x}{9}-\frac{y}{6}\right)+\frac{y}{2}=1.
\end{equation}
要使点与参数 $t$ 无关，需要
\[
\frac{x}{9}-\frac{y}{6}=0,
\qquad
\frac{y}{2}=1.
\]
解得
\begin{equation}
\boxed{(x,y)=(3,2)}.
\end{equation}
所以动直线 $MN$ 恒过定点 $(3,2)$。

从射影观点看，$(3,2)$ 正是固定直线 $AD$ 关于椭圆的极点。这个例题完整体现了视频的命题视角：
\[
Q\text{ 在定直线 }AD\text{ 上运动}
\Longrightarrow
\polar(Q)=MN\text{ 恒过 }\polar(AD).
\]

\section{对高中圆锥曲线复习的提炼}

\subsection*{1. 学生应该真正理解的内容}

\begin{enumerate}
\item 极线公式不是孤立替换，而来自割线上的调和共轭点轨迹；
\item 完全四边形的三个对角点组成自极三角形；
\item La Hire 定理把“点在定直线上”翻译为“对应极线过定点”；
\item 中点、平行、斜率定值等高中条件，可能是调和结构的仿射表现。
\end{enumerate}

\subsection*{2. 考场上可执行的识别顺序}

\begin{enumerate}
\item 看到圆锥曲线上四点和多组对边交点，先补出完全四边形的三个对角点；
\item 判断目标直线是否是某个对角点的极线；
\item 若对应极点在定直线上运动，立即转化为“动极线过该定直线的极点”；
\item 最后用极线方程、韦达定理或点差法写成高中范围内可评分的证明。
\end{enumerate}

\subsection*{3. 教学边界}

射影几何适合解释命题结构和发现结论，但不能替代基本解析运算。面对普通一轮复习学生，建议采用三层结构：
\[
\text{先用坐标法做会}
\quad\Longrightarrow\quad
\text{再用极点极线看懂结构}
\quad\Longrightarrow\quad
\text{最后回到规范书面表达}.
\]

视频提到 2008 年安徽卷、2017 年北京卷、2020 年北京卷、2022 年全国卷和 2023 年全国卷等题目与相关模型有关。本文只记录视频中的这一说法，尚未逐题核对题号和题设；正式编入专题题库前，应再与原题逐一对应。

\section*{一句话总结}

\begin{center}
\fbox{\parbox{0.88\textwidth}{
圆锥曲线中的许多定点、定直线、共线和共点问题，表面上是坐标运算，背后往往是：完全四边形产生调和结构，调和结构产生极点极线，而点线对偶把动点约束转化为定点结论。
}}
\end{center}

\end{document}
